\documentclass[../main]{subfiles}

\begin{document}

\chapter{Partially hyperbolic dynamics}

\section{Partial hyperbolicity}

\begin{definition}[Splitting]
  We say \(TM = E_1\oplus \cdots \oplus E_n\) is a \emph{splitting} of \(TM\) if at each \(x \in M\), \(T_xM = E_1(x)\oplus \cdots \oplus E_n(x)\) and \(\dim E_i(x)\) does not depend on \(x\). We say that the splitting is \emph{continuous}, \emph{smooth}, etc.\ if the subbundle \(E_i \subseteq TM\) is respectively continuous, smooth, etc. (For a definition of subbundles, see~\cite[pp. 264]{leeIntroductionSmoothManifolds2012}.)
\end{definition}

\begin{definition}[Dominated splitting]
  We say a splitting \(TM = E_1\oplus \cdots \oplus E_n\) is \emph{dominated} for a diffeomorphism \(f \colon M \to M\) if:
  \begin{enumerate}
  \item the splitting is \(df\)-\emph{invariant}, that is, \(df_x(E_i(x)) = E_i(f(x))\) for all \(i,x\);
  \item there exists constants \(c > 0\) and \(\lambda \in (0,1)\) such that for all \(1 \le i \le n - 1\), \(x \in M\) \(u \in E_i(x) \setminus \set{0}\), \(v \in E_{i + 1}(x) \setminus \set{0}\) and \(k \ge 0\),
    \begin{equation*}
      \frac{\norm{df^k_x(u)}}{\norm{u}} \le c \lambda^k \frac{\norm{df^k_x(v)}}{\norm{v}}.
    \end{equation*}
  \end{enumerate}
\end{definition}

\begin{definition}
  A subbundle \(E \subseteq TM\) is \emph{uniformly contracted} by \(df\) if there exists \(c > 0\) and \(\lambda \in (0,1)\) such that for all \(x \in M\), \(u \in E(x)\) and \(k \ge 0\),
  \begin{equation*}
    \norm{df^k_x(u)} \le c \lambda^k \norm{u}.
  \end{equation*}
\end{definition}

\begin{definition}
  Let \(f \colon M \to M\) be a diffeomorphism. We say \(f\) is \emph{partially hyperbolic} if there exists a dominated splitting \(TM = E^s \oplus E^c \oplus E^u\) such that \(E^s\) is uniformly contracted by \(df\) and \(E^u\) is uniformly contracted by \(df ^{-1}\).
\end{definition}

\section{Invariance Principle}

This exposition is based on a minicourse given by Maurício Poletti and Yuri Lima on ICTP in 2024, and also on the article~\cite{tahzibiInvariancePrincipleRigidity2018} which the minicourse sourced for the proof.

% \begin{definition}
%   A partially hyperbolic diffeomorphism is \textit{accessible} if between any two points in the manifold, there is a path composed of segments tangent to either \(E^s\) or \(E^u\).
% \end{definition}

\begin{definition}
  A partially hyperbolic diffeomorphism is \textit{dynamically coherent} if there exists invariant foliations \(\fF^{cs}\) and \(\fF^{cu}\) tangent to \(E^c \oplus E^s\) and \(E^c \oplus E^u\), respectively. In this case, we can also define \(\fF^c = \fF^{cs}\vee \fF^{cu}\).
\end{definition}

\begin{definition}
  Consider the quotient map \(\pi \colon X \to X/\mathcal{F}^c\). We say that a partially hyperbolic diffeomorphism \(f \colon X \to X\) has \emph{global unstable holonomies} if for every \(p_1,p_2 \in X/\mathcal{F}^c\) that are in the same leaf of \(\pi(\mathcal{F}^{u})\), there exists an homeomorphism \(h^u_{p_1,p_2}\colon \mathcal{F}^c(p_1) \to \mathcal{F}^c(p_2)\) with the property that \(h^u_{p_1,p_2}(x) \in \mathcal{F}^u(p_1) \cap \mathcal{F}^c(p_2)\). Moreover, for all \(p_1,p_2,p_3 \in X\) in the same leaf of \(\pi(\mathcal{F}^{u})\),
  \begin{enumerate}
  \item \(h^u_{p_1,p_2} \circ h^u_{p_2,p_3} = h^u_{p_1,p_3}\) and
  \item \(h^u_{p_1,p_1} = \Id\).
  \end{enumerate}
  By a slight abuse of notation, if \(p_1,p_2 \in X\), we denote \(h^u_{p_1,p_2} = h^u_{\pi(p_1),\pi(p_2)}\).
\end{definition}

\begin{theorem}\label{thm:invariance-principle}
  Let \(f \colon X \to X\) be a diffeomorphism with compact center leaves and such that the quotient \(f_c \colon X/\mathcal{F}^c \to X/\mathcal{F}^c\) is hyperbolic. Let \(\mu\) be an ergodic invariant measure for \(f\), and also suppose \(f\) has global unstable holonomies. Then, \(\lambda^c(\mu) \le 0\) implies that the disintegration \(\set{\mu^c_x}_{x \in M}\) is invariant by unstable holonomies, meaning that there exists a full measure subset \(E \subseteq X\) in which \((h^u_{p_1,p_2})_*\mu_{p_1}^c = \mu^c_{p_2}\) for all \(p_1,p_2 \in E\) that are in the same \(\mathcal{F}^{cu}\) leaf.
\end{theorem}

\subsection{A trick in a box}

We will begin by explaining the main idea behind the proof. For that, assume we work locally in a leaf of \(\mathcal{F}^{cu}\), in a local coordinate system parametrized by a box \(B = [0,1]^{u} \times [0,1]^{c}\), where \(u = \dim E^u\) and \(c = \dim E^c\). In this box, each leaf of the \(\mathcal{F}^c\) foliation is a ``vertical'' section \(\mathcal{F}^c(x,y) = \set{(x,y) \where y \in [0,1]^{c}}\), and the unstable holonomy maps \(h_{x_1,x_2}^u(x_1,y) = (x_2,y)\) are the identity on \([0,1]^{c}\). Suppose that in this particular leaf, the conditional measure \(\mu^{cu}\) is defined. Also let \(\pi^u(x,y) = x\), \(\pi^c(x,y) = y\), \(\mu^u = \pi^u_*\mu^{cu}\) and \(\mu^c = \pi^c_*\mu^{cu}\).

\begin{figure}[h]
  \centering
  \includegraphics[page=1,scale=0.8]{../figures/invp}
  \caption{Local trivialization of \(\mathcal{F}^{cu}\).}
\end{figure}

We can disintegrate the measure \(\mu^{cu}\) along the foliations \(\mathcal{F}^u\) and \(\mathcal{F}^c\), giving us two families \(\set{\mu^u_p}_{p \in B}\) and \(\set{\mu^c_p}_{p \in B}\). The key fact is that in the box, the projections \(\pi^u_*\mu^u_p\) and \(\pi^c_*\mu^c_p\) uniquely determine \(\mu^u_p\) and \(\mu^c_p\):

\begin{lemma}
  If \(\nu\) is such that \(\nu(\mathcal{F}^c(p)) = 1\), then \(\nu = \delta_{\pi^u(p)} \otimes \pi^c_*\nu\). Similarly, if \(\nu(\mathcal{F}^u(p)) = 1\), then \(\nu = \pi^u_*\nu \otimes \delta_{\pi^c(p)}\).
\end{lemma}

\begin{proof}
  In the first case,
  \begin{align*}
    \nu(A) &= \nu(A \cap \mathcal{F}^c(p))\\
    &= \nu(\set{(x,y) \in A \where x = \pi^u(p)})\\
    &= (\delta_{\pi^u(p)} \otimes \pi^c_*\nu)(A).
  \end{align*}
  The second case is analogous.
\end{proof}

As we remarked before, \(\pi^c \circ h^u_{p_1,p_2} = \Id\). Together, these conditions let us rephrase the invariance principle as the statement that in this box, \(\pi^c_*\mu^c_p = \mu^c\) for \(\mu^{cu}\)-almost every \(p\). By Fubini, this is equivalent to \(\mu^{cu} = \mu^u \otimes \mu^c\) being a product measure, which in turn is equivalent to \(\pi^u_*\mu^u_p = \mu^u\) for \(\mu^{cu}\)-almost every \(p\). Thus, at least locally, we can show the invariance principle by showing that the conditional measures \(\mu^u_p\) all project down to the same measure on the base. 

\subsection{From the box to a leaf}\label{sec:from-box-leaf}

The previous observation is not at all restricted to the box of the last section. We now state the results for an actual ``plaque'' of a \(\mathcal{F}^{cu}\) leaf, which will be the element of a suitable partition.

Denote by \(\pi \colon X \to X/\mathcal{F}^c\) the quotient map, and let \(\set{\mu^c_p}_{p \in X}\) be the disintegration of \(\mu\) along the (measurable) foliation \(\mathcal{F}^c\).
Since \(f_c \colon X/\mathcal{F}^c \to X/\mathcal{F}^c\) is hyperbolic, the quotient also has stable and unstable foliations, which we denote by \(W^s\) and \(W^u\) respectively. Also note that we have \(\mathcal{F}^{cu} = \pi ^{-1}(W^u)\) and \(\mathcal{F}^{cs} = \pi ^{-1}(W^s)\). Let \(\nu = \pi_*\mu\) be the quotient measure on \(X/\mathcal{F}^c\).

We will work with a partition \(\xi_c\) of \(X/\mathcal{F}^c\) whose elements are compact sets and which is finer than the \(W^u\) foliation, that is, \(W^u \prec \xi_c\). By measurability, we have an associated disintegration \(\set{\nu^u_x}_{x \in X/\mathcal{F}^c}\) of \(\nu\).
Denote \(\xi^{cu} = \pi ^{-1}(\xi_c)\), which by construction is finer than the foliation \(\mathcal{F}^u\). Clearly, the elements of \(\xi^{cu}\) are closed, so it is also disintegrates \(\mu\) to a family \(\set{\mu^{cu}_p}_{p \in X}\).
Also let \(\xi^u = \xi^{cu} \vee \mathcal{F}^u\), whose elements are also closed, and gives a disintegration \(\set{\mu^{u}_p}_{p \in X}\).

By the properties of disintegrations, for \(x\) in a full \(\mu\)-measure set \(E_1\), \(\mu^{c}_x(\mathcal{F}^{c}(x)) = \mu^u_x(\xi^u(x)) = 1\) and \(\set{\mu^u_y}_{y \in X}\) is also a disintegration for \(\mu^{cu}_x\) along \(\xi^u\). Since
\begin{equation*}
  \int \mu_x^{cu}(E_1) \;d\mu(x) = \mu(E_1) = 1,
\end{equation*}
\(\mu_x^{cu}(E_1) = 1\) for \(x\) in a full \(\mu\)-measure set \(E_2\). Let \(E_3 \subseteq E_2\) be of full \(\mu\)-measure and such that \(x \in E_3\) implies \(\pi_*\mu^{cu}_x = \nu_{\pi(x)}\).

We now fix \(x \in E_3\), \(\mu^{cu} = \mu^{cu}_x\), \(\nu = \pi_*\mu^{cu}\) and \(P = \xi^{cu}(x)\). We know \(P\) is of the form \(\pi ^{-1}(Q)\) for some \(Q \in \xi_c\). We can parameterize \(P\) by the measurable isomorphism \(\psi \colon Q \times \mathcal{F}(p) \to P\) given by \(\psi(x,q) = h^u_{\pi(p), x}(q)\), whose inverse is \(\psi ^{-1}(q) = (\pi(q),h^u_{\pi(q),\pi(p)}(q))\).

Pushing the measures to the product along the isomorphism, have the same ingredients as in the case of a box. Therefore, we can state a sufficient condition for invariance.

\begin{lemma}\label{lem:proj-implies-invariance}
  For all \(x \in E_3\), if \(\pi_* \mu^u_p = \nu_{\pi(p)}^u\) for \(\mu_x^{cu}\)-almost all \(p\), then there is a full \(\mu_x^{cu}\)-measure set \(E_4(x) \subseteq \xi^{cu}(x)\) where \((h^u_{q_1,q_2})_*\mu^c_{q_1} = \mu^c_{q_2}\) for all \(q_1,q_2 \in E_4(x)\).
\end{lemma}

\begin{proof}
  Restricting to \(x \in E_3\), this is the same argument of the last section, but using Fubini in the product space \(Q \times \mathcal{F}(p)\).
\end{proof}

\subsection{Jensen's inequality}\label{sec:jensens-inequality}

Fix \(x \in X\), and let \(\mathcal{P}\) be a partition of finite entropy for \(\mu^{cu}_x\). We have (see \cref{sec:entropy})
\begin{equation*}
  H_{\mu_x^{cu}}(\mathcal{P}) = \sum_{P \in \mathcal{P}} I(\mu_x^{cu}(P)),
\end{equation*}
and also the conditional entropy
\begin{equation*}
  H_{\mu_x^{cu}}(\mathcal{P} \mid \mathcal{F}^u)
  = \int H_{\mu^{u}_y}(\mathcal{P}) \;d\mu_x^{cu}(y)
  = \int \sum_{P \in \mathcal{P}} I(\mu_y^{u}(P)) \;d\mu_x^{cu}(y).
\end{equation*}

\begin{figure}[h]
  \centering
  \includegraphics[page=2,scale=0.8]{../figures/invp}
  \caption{Schematic of a partition \(\mathcal{P}\) and the measures \(\mu^u_y\).}
\end{figure}

\begin{proposition}[Jensen's Inequality]\label{prop:jensen}
  Suppose \(\varphi \colon \RR \to \RR\) is a concave function, and let \(\psi \colon \Omega \to \RR\) be a measurable function for a probability space \((\Omega,\mathcal{B},\Prob)\). Then
  \begin{equation*}
    \int \varphi \circ \psi \;d\Prob \le \varphi\rpar[\bigg]{\int \psi \;d\Prob}.
  \end{equation*}
  Moreover, if \(\varphi\) is strictly concave, the equality holds if, and only if, \(\psi\) is constant almost everywhere.
\end{proposition}
% \begin{proposition}[Jensen's Inequality]\label{prop:jensen}
%   Suppose \(E\) is a real Banach space, \(\varphi \colon E \to \RR\) is continuous and concave, and \(\psi \colon \Omega \to E\) is a Bochner integrable function defined on a probability space \((\Omega,\mathcal{B},\Prob)\). Then
%   \begin{equation*}
%     \int \varphi \circ \psi \;d\Prob \le \varphi\rpar[\bigg]{\int \psi \;d\Prob}.
%   \end{equation*}
%   Moreover, if \(\varphi\) is strictly concave, the equality holds if, and only if, \(\psi\) is constant almost everywhere.
% \end{proposition}

Applying \cref{prop:jensen} to \(\Prob = \mu_x^{cu}\), \(\varphi(x) = I(x) =  - x\log x\) (which is strictly concave) and \(\psi_P(y) = \mu^u_p(P)\), for each \(P \in \mathcal{P}\) we have
\begin{equation*}
  \int I(\mu_p^u(P))\;d\mu_x^{cu}(p)\le I\rpar[\bigg]{\int \mu_p^u(P)\;d\mu_x^{cu}(p) } = I(\mu_x^{cu}(P)),
\end{equation*}
so summing over \(P \in \mathcal{P}\),
\begin{equation}\label{eq:5}
  H_{\mu_x^{cu}}(\mathcal{P}\mid \mathcal{F}^u)\le H_{\mu_x^{cu}}(\mathcal{P})
\end{equation}
with equality if, and only if, \(p \mapsto \mu^u_p(P)\) is constant for \(\mu_x^{cu}\)-every \(p\) and all \(P\).

Now, suppose \(\mathcal{P}\) is of the form \(\mathcal{P} = \pi ^{-1}(\xi_c)\), for some partition \(\xi_c\) of \(X/\mathcal{F}^c\). If \(H_{\mu_x^{cu}}(\mathcal{P}\mid \mathcal{F}^u) = H_{\mu_x^{cu}}(\mathcal{P})\), then as we noticed, \(p \mapsto \mu^u_p(P)\) is constant for \(\mu_x^{cu}\)-every \(p\) and all \(P \in \mathcal{P}\). Suppose additionally that \(x \in E_3\) as defined in the last section. Then, for every \(A \in \xi_c\),

\begin{equation*}
  \int \pi_*\mu_p^{u}(A) \;d\mu_x^{cu}(p) =
  \int \mu_p^{u}(\pi ^{-1}(A)) \;d\mu_x^{cu}(p) =
  \mu_x^{cu}(\pi ^{-1}(A)) = \nu_{\pi(x)}(A),
\end{equation*}
so \(\pi_*\mu_p^{u}(A) = \nu_{\pi(x)}(A)\) for \(\mu^{cu}_x\)-almost all \(p\) and all \(A \in \xi_c\).
In other words, we have shown the following lemma.

\begin{lemma}\label{lem:almost-there}
  Let \(x \in E_3\) and let \(\mathcal{P} = \pi ^{-1}(\xi_c)\) be a partition with finite entropy for \(\mu^{cu}_x\). If \(H_{\mu_x^{cu}}(\mathcal{P}\mid \mathcal{F}^u) = H_{\mu_x^{cu}}(\mathcal{P})\), then for \(\mu^{cu}_x\)-almost all \(p\), \(\pi_*\mu^u_p\) and \(\nu^u_p\) coincide in the \(\sigma\)-algebra \(\sigma(\xi_c)\).
\end{lemma}

\begin{proof}
  See above.
\end{proof}

\begin{lemma}\label{cor:gen-implies-inv}
  Fix \(\xi_c\) and suppose that \(\sigma(\bigcup_n f_c^{-n}(\xi_c))\) is finer than the Borel algebra of \(X/\mathcal{F}^c\). Let \(\mathcal{P}_n = \pi ^{-1}(f_c^{-n}(\xi_c))\). For all \(x \in E_3\), if \(H_{\mu_x^{cu}}(\mathcal{P}_n\mid \mathcal{F}^u) = H_{\mu_x^{cu}}(\mathcal{P}_n)\) for all \(n\), then there exists a full \(\mu_x^{cu}\)-measure set \(E_4(x)\), such that for all \(p_1, p_2 \in E_4(x)\), \((h^u_{p_1,p_2})_*\mu^c_{p_1} = \mu^c_{p_2}\).
\end{lemma}

\begin{proof}
  Let \(x \in E_3\). By \cref{lem:almost-there}, for all \(n\), \(\pi_*\mu^u_p\) and \(\nu^u_p\) coincide in the \(\sigma\)-algebra \(\sigma(f_c^{-n}(\xi_c))\) for \(\mu^{cu}_x\)-almost all \(p\). Since \(\bigcup_n f_c^{-n}(\xi_c)\) is finer than the Borel algebra, \(\pi_*\mu^u_p = \nu^u_p\) for \(\mu^{cu}_x\)-almost all \(p\). The result follows by applying \cref{lem:proj-implies-invariance}, which gives the set \(E_4(x)\).
\end{proof}

\begin{lemma}\label{cor:gen-implies-inv-ae}
  Suppose that \(\sigma(\bigcup_n f_c^{-n}(\xi_c))\) is finer than the Borel algebra of \(X/\mathcal{F}^c\), and in the notation of \cref{cor:gen-implies-inv},  \(H_{\mu_x^{cu}}(\mathcal{P}_n\mid \mathcal{F}^u) = H_{\mu_x^{cu}}(\mathcal{P}_n)\) for all \(n\) and \(\mu\)-almost every \(x\). Then, there is a full \(\mu\)-measure set \(E_5 \subseteq X\) such that for all \(x,y \in E_5\), if \(y \in \xi^{cu}(x)\), then \((h^u_{x,y})_*\mu^c_x = \mu^c_y\).
\end{lemma}

\begin{proof}
  Let \(E \subseteq E_3\) be a full-measure subset of \(E_3\) where \(H_{\mu_x^{cu}}(\mathcal{P}_n\mid \mathcal{F}^u) = H_{\mu_x^{cu}}(\mathcal{P}_n)\) holds for all \(n\). We explicitly define
  \begin{equation*}
    E_5 = \bigcup_{x \in E}E_4(x).
  \end{equation*}
  This set has full measure because
  \begin{equation*}
    \mu(E_5) \ge \int_{E} \mu^{cu}_x(E_4(x)) \;d\mu(x) \ge \mu(E) = 1.
  \end{equation*}
  Now, if \(x,y \in E_5\), then by definition there exists \(p \in E\) such that \(x \in E_4(p) \subseteq \xi^{cu}(p)\). Therefore, if \(y \in \xi^{cu}(x)\), then \(y \in \xi^{cu}(p)\) and \(y \in E_4(p)\). By definition of \(E_4(p)\) and \cref{cor:gen-implies-inv}, we know \((h^u_{x,y})_*\mu^c_x = \mu^c_y\).
\end{proof}

\subsection{Patching the argument together}

We have shown a sufficient condition for the invariance principle to hold for a pair of points inside the \emph{same} \(cu\)-plaque, that is, an element of \(\xi^{cu}\). But generally, even if \(x\) and \(y\) are in the same leaf of \(\mathcal{F}^{cu}\), they need not be in the same element of \(\xi^{cu}\). So we need a small adaptation to show \((h^u_{x,y})_*\mu^c_x = \mu^c_y\) in this case. For that, we assume an extra condition on \(\xi_c\):

\begin{hypothesis}\label{hip:neighb}
  For \(\mu\)-almost every \(x\), \(\xi^u(x) = \xi^{cu}(x) \cap \mathcal{F}^u(x)\) is a neighborhood of \(x\) in \(\mathcal{F}^u(x)\).
\end{hypothesis}

\begin{lemma}
  There exists \(\delta > 0\) such that the set
  \begin{equation*}
    S = S_{\delta} = \set{x \in X \where d^u(x, X \setminus \xi^{u}(x)) > \delta}
  \end{equation*}
  has positive measure.
\end{lemma}
\begin{proof}
  By a simple application of \cref{hip:neighb} and the monotone convergence theorem, \(\lim_{\delta \to 0}      \mu(S_{\delta})\to 1\). Choose some \(\delta > 0\) for which \(\mu(S_{\delta})\) is positive.
\end{proof}

We finally arrive at the global version.

\begin{lemma}\label{cor:cond-invprin}
  Suppose that \(\sigma(\bigcup_n f_c^{-n}(\xi_c))\) is finer than the Borel algebra of \(X/\mathcal{F}^c\), and in the notation of \cref{cor:gen-implies-inv},  \(H_{\mu_x^{cu}}(\mathcal{P}_n\mid \mathcal{F}^u) = H_{\mu_x^{cu}}(\mathcal{P}_n)\) for all \(n\) and \(\mu\)-almost every \(x\). Then, there is a full \(\mu\)-measure set \(E \subseteq X\) such that for all \(x,y \in E\), if \(y \in \mathcal{F}^{u}(x)\), then \((h^u_{x,y})_*\mu^c_x = \mu^c_y\).
\end{lemma}

\begin{proof}
  Choose a full-measure \(E\) such that
  \begin{enumerate}
  \item \(\mu^c_{f^k(x)} = f^k_*\mu^c_x\) for all \(k \in \ZZ\), for all \(x \in E\) (by invariance of \(\mu\));
  \item \(f^{-n}(x) \in S\) for infinitely many \(n > 0\), for all \(x \in E\) (by ergodicity);
  \item \(E \subseteq \bigcap_{k \in \NN} f^{k}(E_5)\), where \(E_5\) is given by \cref{cor:gen-implies-inv-ae}.
  \end{enumerate}
  Let \(x,y \in E\) such that \(y \in \mathcal{F}^u(x)\). We know \(d^u(f^{-n}(x),f^{-n}(y)) \to 0\), so there exists some \(N\) for which \(d^u(f^{-N}(x),f^{-N}(y)) < \delta\) and \(f^{-N}(x) \in S\). By definition of \(S\), this implies \(f^{-N}(y) \in \xi^{cu}(f^{-N}(x))\). Since \(x, y \in f^N(E_5)\), \cref{cor:gen-implies-inv-ae} implies
  \[f^{-N}_*(h^u_{x,y})_*\mu^c_x = (h^u_{f^{-N}(x),f^{-N}(y)})_*\mu^c_{f^{-N}(x)} = \mu^c_{f^{-N}(y)} = f^{-N}_*\mu^c_y.\]
  Now, apply \(f^N_*\) to both sides to conclude \((h^u_{x,y})_*\mu^c_x = \mu^c_y\).
\end{proof}

\subsection{Adapted partitions and unstable entropy}

We now define and construct special kind of partition refining the unstable foliation, and an associated notion of entropy, that will satisfy the hypothesis of \cref{cor:cond-invprin}.

\begin{definition}[Adapted partition]
  We say a measurable partition \(\xi\) is \(\mu\)-adapted to a foliation \(\mathcal{F}\) if \(\xi \prec f^{-1}(\xi)\) and there exists \(\varepsilon > 0\) such that for \(\mu\)-\almostevery \(x \in M\),
  \begin{enumerate}
  \item\label{item:2} \(\xi(x) \subseteq \mathcal{F}(x)\cap B_{\varepsilon}(x)\) and
  \item\label{item:4} \(\xi(x)\) is a neighborhood of \(x\) inside \(\mathcal{F}(x)\).
  \end{enumerate}
\end{definition}

\begin{figure}[h]
  \centering
  \includegraphics[page=3,scale=0.8]{../figures/invp}
  \caption{A partition \(\xi\) adapted to \(\mathcal{F}^u\).}
\end{figure}

\begin{definition}[Unstable entropy]
  Let \(\xi\) be a partition \(\mu\)-adapted to the unstable foliation \(\mathcal{F}^u\). The \emph{unstable entropy} is defined as
  \begin{equation*}
    h_{\mu}^u(f) =
    h_{\mu}(f, \xi) =
    H_{\mu}(f ^{-1}(\xi) \mid \xi) =
    \int H_{\mu_x^u}(f^{-1}(\xi))\;d\mu(x),
 \end{equation*}
 where \(\set{\mu^u_x}_{x \in X}\) is the decomposition of \(\mu\) with respect to \(\xi\).
\end{definition}

% \begin{remark}
%   Notice that the partitions \(f^{-1}(\xi)\) have finite entropy for \(\mu_x^u\) for \(\mu\)-almost every \(x\), so the definition makes sense. Indeed, we know
% \end{remark}

The following proof is expanded from~\cite[lemma 3.1.2]{ledrappierMetricEntropyDiffeomorphisms1985}.
\begin{lemma}
  The unstable entropy does not depend on the choice of adapted partition \(\xi\).
\end{lemma}

\begin{proof}
  Let \(\xi_1\) and \(\xi_2\) be two partitions adapted for the unstable foliation. It suffices to show that \(h_{\mu}(f, \xi_1 \vee \xi_2) = h_{\mu}(f,\xi_1)\). Let  \(n \ge 0\), and notice that by monotonicity,
  \begin{equation*}
    h_{\mu}(f, \xi_1 \vee \xi_2)\ge h_{\mu}(f, \xi_1 \vee f^n(\xi_2)).
  \end{equation*}
  On the other hand,
  \begin{equation*}
     h_{\mu}(f, \xi_1 \vee \xi_2) \le h_{\mu}(f, f^{-n}(\xi_1) \vee \xi_2) = h_{\mu}(f, \xi_1 \vee f^n(\xi_2)).
  \end{equation*}
  So
  \begin{align*}
    h_{\mu}(f, \xi_1 \vee \xi_2) &= h_{\mu}(f, \xi_1 \vee f^n(\xi_2))\\
    &= H_{\mu}(\xi_1 \vee f^n(\xi_2) \mid f(\xi_1) \vee f^{n + 1}(\xi_2))\\
    &= H_{\mu}(\xi_1 \mid f(\xi_1) \vee f^{n + 1}(\xi_2))
    +  H_{\mu}(\xi_2 \mid f^{-n + 1}(\xi_1) \vee f(\xi_2)).
  \end{align*}
  Therefore,
  Now, we let \(n \to \infty\). Since \(\xi_2 \prec \bigvee_n f^{-n + 1}(\xi_1)\), the second term goes to zero. For the first term, define the set \(D_n = \set{x \in X \where f(\xi_1)(x) \subseteq f^{n + 1}(\xi_2)(x)}\).

  \begin{claim}
    \(\mu(D_n) \to 1\).
  \end{claim}
  To show the claim, we use that \(f\) is measure-preserving, so
  \begin{align*}
    \mu(D_n) &= \mu(\set{x \in X \where f(\xi_1)(f^{n + 1}(x)) \subseteq f^{n + 1}(\xi_2)(f^{n + 1}(x))})\\
     &= \mu(\set{x \in X \where f^{-n}(\xi_1)(x) \subseteq \xi_2(x)}).
  \end{align*}
  Since for \(\mu\)-almost every \(x \in X\), \(\xi_2(x)\) contains a neighborhood in the unstable leaf and \(\diam f^{-n}(\xi_1)(x) \to 0\), the set \(f^{-n}(\xi_1)(x)\) is eventually contained in \(\xi_2(x)\) for \almostevery \(x\). By the dominated convergence theorem, this shows \(\mu(D_n) \to 1\).
  
  Therefore,
  \begin{align*}
    H_{\mu}(\xi_1 \mid f(\xi_1) \vee f^{n + 1}(\xi_2)) &\ge
    H_{\mu \llcorner D_n}(\xi_1 \mid f(\xi_1) \vee f^{n + 1}(\xi_2))\\ &= 
    H_{\mu \llcorner D_n}(\xi_1 \mid f(\xi_1)) \to H_{\mu}(\xi_1 \mid f(\xi_1)).
  \end{align*}
  Since \(H_{\mu}(\xi_1 \mid f(\xi_1) \vee f^{n + 1}(\xi_2)) \le H_{\mu}(\xi_1 \mid f(\xi_1))\), we have \(h_{\mu}(f, \xi_1 \vee \xi_2) = h_{\mu}(f,\xi_1)\).
\end{proof}

Let's construct this adapted partition \(\xi\). Since \(f_c \colon X/\mathcal{F}^c \to X/\mathcal{F}^c\) is hyperbolic for \(\nu\), it admits a Markov partition whose boundary is a null \(\nu\)-measure set (see, for instance,~\cite[Theorem 3.12.]{bowenEquilibriumStatesErgodic2008}). For a choice of small \(\varepsilon > 0\), a Markov partition is defined as follows.

\begin{definition}
  We say a set \(R \subseteq X/\mathcal{F}^c\) is a \emph{proper rectangle} if \(R\) is closed, \(R = \closure{\interior{R}}\) and \([x,y] \in R\) whenever \(x, y \in R\). The bracket \([x,y]\) is defined as the unique point in \(W^s_{\varepsilon}(x)\cap W^u_{\varepsilon}(y)\) for all \(x,y\) such that \(d(x,y) < \delta\) for some \(0 < \delta < \varepsilon\) (see~\cite[3.3. Canonical Coordinates]{bowenEquilibriumStatesErgodic2008}). For that to make sense, we also require that \(\diam R < \delta\).
\end{definition}

\begin{definition}
  A \emph{Markov partition} is a finite partition \(\mathcal{R} = \set{R_1,\dots,R_m}\) of \(X/\mathcal{F}^c\) by proper rectangles such that for all \(x \in X\) with \(x, f(x) \notin \bigcup_{i = 1}^m \partial R_i\),
  \begin{enumerate}
  \item\label{item:6} \(W^u(f(x)) \cap \mathcal{R}(f(x)) \subseteq f(W^u(x)\cap \mathcal{R}(x))\) and
  \item \(W^s(f(x)) \cap \mathcal{R}(f(x)) \supseteq f(W^s(x)\cap \mathcal{R}(x))\).
  \end{enumerate}
\end{definition}

Given a Markov partition \(\mathcal{R}\), we define \(\xi_c(x) = W^u(x) \cap \mathcal{R}(x)\). Now, construct \(\xi^{cu}\) and \(\xi^u\) as in \cref{sec:from-box-leaf}. We claim \(\xi = \xi^{u}\) is \(\mu\)-adapted to \(\mathcal{F}^u\). Indeed, by \cref{item:6} we have \(\xi^{u} \prec f^{-1}(\xi^{u})\), and outside the zero-measure set \(\pi ^{-1}\rpar*{\bigcup_{i = 1}^m \partial R_i}\), we have \[\xi^{u}(x) = \mathcal{F}^u(x) \cap \pi^{-1}(W^u(\pi(x)) \cap \mathcal{R}(\pi(x))),\]
so \(\xi^{u}(x)\) is a neighborhood of \(x\) in \(\mathcal{F}^u(x)\) of small diameter because \(\pi(x) \in \interior \mathcal{R}(\pi(x))\).

With the partition \(\xi = \xi^u\) in hand, we finish the proof of \cref{thm:invariance-principle} with a series of lemmas.

\begin{lemma}\label{lem:entropy-eq}
  For \(\mu\)-almost every \(x\),
  \begin{equation*}
    H_{\mu_{x}^u}(f^{-1}(\xi)) = H_{\mu_{x}^u}(f^{-1}(\xi^{cu}))
  \end{equation*}
\end{lemma}

\begin{proof}
  Since \(\xi \prec f^{-1}(\xi)\) and by definition, the elements of \(f^{-1}(\xi)\) have the form
  \[f^{-1}(\xi)(y) = f^{-1}(\xi^{cu})(y) \cap \mathcal{F}^u(y) = f^{-1}(\xi^{cu})(y) \cap \xi(y).\]
  By definition of a disintegration, for \(\mu\)-almost every \(x\) and \(\mu_x^u\)-almost every \(y\), \(\mu_y^u(\xi(y)) = 1\), so
  \begin{align*}
    H_{\mu_{x}^u}(f^{-1}(\xi))
    &= - \int \log \mu(f^{-1}(\xi)(y)) \;d\mu_x^u(y)\\
    &= - \int \log \mu(f^{-1}(\xi^{cu})(y) \cap \xi(y)) \;d\mu_x^u(y)\\
    &= - \int \log \mu(f^{-1}(\xi^{cu})(y)) \;d\mu_x^u(y) = H_{\mu_{x}^u}(f^{-1}(\xi^{cu})).
  \end{align*}
\end{proof}

\begin{figure}[h]
  \centering
  \includegraphics[page=4,scale=0.8]{../figures/invp}
  \caption{The identity \(f^{-1}(\xi)(x) = f^{-1}(\xi^{cu})(x) \cap \xi(x)\).}
\end{figure}

\begin{proof}[Proof of \cref{thm:invariance-principle}]
We have the inequality
\begin{align*}
  h_{\mu}^u(f)
  = \int H_{\mu_{x}^u}(f^{-1}(\xi))\;d\mu(x)
  &= \int H_{\mu_{x}^u}(f^{-1}(\xi^{cu}))\;d\mu(x)\\
  &= \int H_{\mu_{x}^{cu}}(f^{-1}(\xi^{cu}) \mid \mathcal{F}^u)\;d\mu(x)\\
  &\le \int H_{\mu^{cu}_x}(f^{-1}(\xi^{cu}))\;d\mu(x)\\
  &=  \int H_{\nu^{u}_{\pi(x)}}(f_c^{-1}(\xi_c))\;d\mu(x)\\
  &=  \int H_{\nu^{u}_{y}}(f_c^{-1}(\xi_c))\;d\nu(y)\\
  &= h^u_{\nu}(f_c) = h_{\nu}(f_c)
\end{align*}
with equality if, and only if, \(H_{\mu_{x}^{cu}}(f^{-1}(\xi^{cu}) \mid \mathcal{F}^u) = H_{\mu^{cu}_x}(f^{-1}(\xi^{cu}))\) for \(\mu\)-almost every \(x\).
In the first line, we applied \cref{lem:entropy-eq}, and in the third line, we used  the \cref{eq:5} of \cref{sec:jensens-inequality}. The inequality above equally applies to \(f = f^n\), and since the diameter of \(f^{-n}(\xi_c)\) shrinks to zero, \(\sigma\rpar*{\bigcup_nf^{-n}(\xi_c)}\) is finer than the Borel algebra.

Now, since \((f^n_c, \nu)\) is a factor of \((f^n,\mu)\), we have \(h_{\nu}(f_c^n) \le h_{\mu}(f^n)\). If \(\lambda^c(\mu) \le 0\), then the Ledrappier-Young formula tells us that \(h_{\mu}^u(f^n) = n h_{\mu}^u(f) = n h_{\mu}(f) = h_{\mu}(f^n)\), so we have equality in the equation above. By \cref{cor:cond-invprin}, the proof of \cref{thm:invariance-principle} is complete.
\end{proof}

\section{Applications to Bonifant-Milnor}

\todo{todo}

% We begin the proof of \cref{thm:invariance-principle-skew-products} by examining a trick for establishing the invariance of the conditional measures by the unstable holonomy. Let 


\end{document}
